\section{The actual finite total cup and original section pairing}

\begin{definition}[The signed finite total cup]
\label{def:di-cech-total-cup}
\uses{def:native-cech-ordered-cup}
\leanok
\lean{EnvelopeProperty.DI.CechTotalCup.totalCupComponent,EnvelopeProperty.DI.CechTotalCup.totalCup}
Let $Y=X^{(p)}$, let $K$ be the original pushed de Rham complex
of a smooth morphism of pure relative dimension $n$, and let $U$ be any
family of opens of $Y$. Let $T=\operatorname{Tot}\check C(U,K)$ be the
original first-quadrant total complex. For $q+j=a$ and $r+j'=b$, take
the original projections from $T^a,T^b$ to these two bidegrees, their
original tensor map, the ordered Alexander--Whitney wedge product,
and the original inclusion of bidegree $(r+q,j+j')$ in $T^{a+b}$.
Multiply this composite by $(-1)^{jr}$ and sum over the two finite
antidiagonals. This defines an actual sheaf morphism
$T^a\otimes T^b\to T^{a+b}$. Neither a cover nor a finite indexing
set is required: each total degree has only finitely many bidegrees.
\end{definition}
\begin{proof}
The published wedge has form degree $r+q$, and the ordered cup has
Cech degree $j+j'$. The antidiagonal equalities imply
$(r+q)+(j+j')=a+b$, which is the equality used in the original total
inclusion. Every summand has the same tensor source and total target,
so the two finite sums exist in the original additive module category.
\end{proof}

\begin{definition}[The whole original bilinear section pairing]
\label{def:di-cech-total-cup-section}
\uses{def:di-cech-total-cup}
\leanok
\lean{EnvelopeProperty.DI.CechTotalCup.totalCupSection}
For every open $W\subseteq Y$, compose the original tensor-section
bilinear map on $\Gamma(W,T^a),\Gamma(W,T^b)$ with the section map
of the actual total cup morphism. This is a bilinear pairing on every
pair of original sections, with values in $\Gamma(W,T^{a+b})$.
\end{definition}
\begin{proof}
The original tensor-section pairing is bilinear over
$\Gamma(W,\mathcal O_Y)$, and the section map of a sheaf-module
morphism is linear over that same ring. Their composition is bilinear.
\end{proof}

\begin{theorem}[Original tensor evaluation and restriction]
\label{thm:di-cech-total-cup-section-restriction}
\uses{def:di-cech-total-cup-section}
\leanok
\lean{EnvelopeProperty.DI.CechTotalCup.totalCupSection_tensorSection,EnvelopeProperty.DI.CechTotalCup.totalCupSection_restrict}
The entire section pairing is the actual total cup applied to the
original tensor section. It commutes with restriction along every
open inclusion, for every pair of whole original sections.
\end{theorem}
\begin{proof}
The tensor evaluation identity is the definition of composition of the
original bilinear map with the section map. The original tensor section
commutes with restriction, by naturality of the original sheafification
unit. Naturality of the actual total cup sheaf morphism then identifies
restriction of its value with its value on the two restricted sections.
\end{proof}

The differential Leibniz identity and its closedness consequence are
separate subsequent results, not hypotheses or claims of this construction.
